% FILE: 00_project_identity.tex
% Project: research-prompt-manufactory
% Thesis Role: prompt-manufacturing-and-subagent-command-surface

\section{Project Identity: Prompt Manufactory}
\label{sec:research-prompt-manufactory-identity}

\subsection{Purpose}
\label{sec:research-prompt-manufactory-identity-purpose}

The Prompt Manufactory is the command-surface manufacturing layer of the Process Intelligence
research program. Its purpose is to manufacture typed, receipted, ontology-grounded research
warrants for downstream subagents — with zero hand-written prompt text permitted. Every research
warrant emitted by the Manufactory carries a \texttt{pm:derivedFrom} triple attesting its
provenance to a specific RDF instance in graph law, enforced by a SHACL shape constraint that
rejects any \texttt{pm:RenderedPrompt} lacking this triple.

The doctrine governing the Manufactory is stated explicitly in the checkpoint record
\texttt{GGEN\_PROMPT\_MANUFACTORY\_PARTIAL\_001.md}: \textit{``The prompt is no longer speech.
It is a receipted production order emitted from graph law.''} This doctrine inverts the
conventional assumption that prompts are human-authored natural-language instructions. Instead,
a prompt in this system is a manufactured artifact with a traceable derivation chain, a type
classification, and a receipt entry in the \texttt{prompt-receipt-ledger.md}.

The Manufactory serves seven research programs encoded as seed instances in
\texttt{research-program-law.ttl}: PI, GGEN\_ECOSYSTEM, GGEN\_OTEL, ZOEAPP, EXPO\_SUPABASE,
CLAUDE\_WORKFLOW, and WASM4PM\_REMEDIATE. For each program, the Manufactory emits workflow
warrants, subagent role warrants, checkpoint warrants, skill warrants, and hook policy
documents --- all manufactured from the same graph.

\subsection{Repository Surface}
\label{sec:research-prompt-manufactory-identity-surface}

The Prompt Manufactory occupies the path
\texttt{/Users/sac/process-intelligence/research/prompt-manufactory/} within the
\texttt{process-intelligence} research foundry. Its internal layout enforces strict separation
of concerns across four layers:

\begin{itemize}
  \item \textbf{Ontology layer} (\texttt{ggen/ontology/}): Eight Turtle files (\texttt{.ttl})
    encoding all law. Each file governs a distinct domain: prompt-manufactory (core class
    definitions), research-program-law (7 seed program instances), workflow-law (INTEL and
    REMEDIATE workflow structures with 8 phases), subagent-role-law (15 bounded inspection
    station definitions), skill-law (6 standard-work capabilities), hook-law (6 Andon gate
    policies), checkpoint-law (ALIVE/PARTIAL verdict structure with 10 gate definitions), and
    forbidden-collapse-law (3 absolute bans plus 22 legacy \texttt{.ggen} file classification
    instances).
  \item \textbf{Query layer} (\texttt{ggen/queries/}): SPARQL SELECT queries that extract
    structured data from the RDF graph for each warrant type. The core query
    \texttt{select-workflow-prompts.rq} joins \texttt{pm:ResearchProgram},
    \texttt{pm:Workflow}, \texttt{pm:Phase}, and \texttt{pm:SubagentRole} across four
    graph hops.
  \item \textbf{Template layer} (\texttt{ggen/templates/}): Tera templates (\texttt{.md.tera})
    that render SPARQL result bindings into Markdown prompt artifacts. Eight template
    specifications are wired in \texttt{ggen.toml}; one proof specimen
    (\texttt{workflow-prompt.md.tera}) is operational.
  \item \textbf{Emission layer} (\texttt{emitted/}): The manufactured artifacts directory,
    organized into \texttt{prompts/workflows/}, \texttt{prompts/subagents/},
    \texttt{prompts/checkpoints/}, \texttt{prompts/skills/}, \texttt{prompts/hooks/},
    \texttt{indexes/}, and \texttt{manifests/}.
\end{itemize}

The manufacturing engine is ggen v26.5.21, configured in \texttt{ggen.toml} with 6 active
generation rules. The engine's sync state is recorded in
\texttt{ggen/.ggen/sync-state.json} with ontology hash
\texttt{01b52b815102f2839b5c852cadf9edb119592635bd8d8436b9574e0aea6bb054} and
sync timestamp \texttt{2026-06-01T23:08:25Z}.

\subsection{Primary Research Role}
\label{sec:research-prompt-manufactory-identity-role}

Within the dissertation, the Prompt Manufactory fills the thesis role of
\textbf{prompt-manufacturing-and-subagent-command-surface}. It is the layer that converts
graph-encoded research law into actionable, typed, receipted instructions that subagents
can execute. This role sits between the ontology layer (where law is declared) and the
execution layer (where subagents act), acting as the manufacturing bridge.

The Manufactory demonstrates that the command surface of a multi-agent research program
need not be hand-authored. If law can be encoded as RDF, and if manufacturing rules can
be specified as SPARQL-plus-Tera pipelines, then every subagent instruction becomes a
manufactured artifact with provenance, type, and receipt --- properties that hand-written
prompts cannot possess.

\subsection{Key Artifacts}
\label{sec:research-prompt-manufactory-identity-artifacts}

The key artifacts produced by or constitutive of the Prompt Manufactory are:

\begin{itemize}
  \item \texttt{ggen/ontology/prompt-manufactory.ttl} --- 15 core classes, 8 object
    properties, 7 data properties, SHACL shapes; the foundational schema.
  \item \texttt{ggen/ggen.toml} --- 6 generation rules wiring SPARQL queries to Tera
    templates; the manufacturing pipeline specification.
  \item \texttt{emitted/manifests/manufacturing-manifest-20260601.yaml} --- Records ggen
    v26.5.21 executing 6 rules, emitting 41 artifacts across 5 types, with overall gate
    PASS and receipts-equal-artifacts PASS.
  \item \texttt{emitted/indexes/prompt-receipt-ledger.md} --- The manufacturing receipt
    ledger documenting the provenance chain for each emitted warrant.
  \item \texttt{emitted/prompts/workflows/PI\_RESEARCH\_PROGRAM\_INTEL\_001.md} --- The
    proof-specimen workflow warrant, demonstrating end-to-end manufacture from
    \texttt{research-program-law.ttl} to receipted Markdown.
  \item \texttt{checkpoints/GGEN\_PROMPT\_MANUFACTORY\_PARTIAL\_001.md} --- The issued
    PARTIAL checkpoint with 9/11 audit gates passing, 2 pending, 0 blocking.
\end{itemize}
